\begin{slide}
\begin{center}\Title{Our main definitions}\end{center}

\begin{itemize}

\IncPart{1}{\item[]
In this setting $T\Subs{W}{V}$ replaces the set of paths $V$ with $W$ in $T$
and $\IS\LabP{p}\LabP{T} \bydef \ExpSubset{\IS\LabP{p}\LabP{q}}{q \in T}$
for $p$ and $q$ arbitrary segments of paths.
}

\IncPart{1}{\item[]
\stress{Definition (subtree).}
Let $\IS\Subtree{T}(p) \bydef \ExpSubset{q}{\IS\LabP{p}\LabP{q} \in T}$. 
}

\IncPart{2}{\item[]
\stress{Balanced segments}
are a key notion in reduction at a distance.
}

\IncPart{2}{\item[]
\stress{Definition (balanced segment).}
$\Empty$ (the empty path segment) is balanced,
if $b$ and $c$ are balanced so are:
$\IS\LabA\LabP{b}\LabL$,
$\IS\LabD{h}\LabP{b}$,
$\IS\LabP{b}\LabD{h}$,
$\IS\LabP{b}\LabP{c}$.
}

\IncPart{3}{\item[]
When $\IS\LabP{p}\LabL\LabP{q}\LabP{(n+1)} \in T$ and $n \in \N$
the occurrence $n+1$ \stress{refers to the displayed} $\IS\LabL$
when $\Depth{q} = n$, with $\Depth{q}$ aware of delayed lifts.
}

\IncPart{4}{\item[]
\stress{Definition (depth).}
$\Depth{p} \in \N$ is given by:
$\Depth{\Empty} \bydef 0$,
$\Depth{\IS\LabP{p}\LabA} \bydef \Depth{\IS\LabP{p}\LabS} \bydef \Depth{p}$,
$\Depth{\IS\LabP{p}\LabL} \bydef \Depth{p}+1$,
if $\Depth{p} \ge h$ then $\Depth{\IS\LabP{p}\LabD{h}} \bydef \Depth{p}-h$.
}

\IncPart{5}{\item[]
\stress{Definition (ef-reduction).}
If $r = \IS\LabP{p}\LabA\LabP{b}\LabL\LabP{q}\LabP{(n+1)} \in T$,
$b$ is balanced and $\Depth{q} = n \in \N$ then
$T \EFStep{r} T\Subs{\IS\LabP{r}\Subtree{T}(\IS\LabP{p}\LabS)}{r}$.
}

\IncPart{6}{\item[]
\stress{Key observation.}
The delayed lifts in the segment $b$
are never applied to any variable occurrence
and thus we can rule them out.
}

\end{itemize}

\end{slide}
